\begin{table}[h]
\centering
\footnotesize
\begin{tabular}{p{.9\columnwidth}}
\toprule
\textbf{Prompt: LLM Judge} \\
\midrule
\end{tabular}
\begin{lstlisting}[style=promptstyle]
###Task Description:
An instruction (might include an Input inside it), a response to evaluate, and a score rubric representing an evaluation criteria are given.
1. Write a detailed feedback that assesses the quality of the response strictly based on the given score rubric, not evaluating in general.
2. After writing a feedback, write a score that is an integer between 1 and 5. You should refer to the score rubric.
3. The output format should look as follows: "Feedback: (write a feedback for criteria) [RESULT] (an integer number between 1 and 5)"
4. Please do not generate any other opening, closing, or explanations.

###Instruction:
You are evaluating a summary from the perspective of a specific person. Consider their background and information needs when assessing quality.

Annotator profile:
- Role: {role}
- Domain: {domain}
- Information needs: {info_needs}
- Query: {query}

Evaluate the quality of the following summary on the dimension of {dimension}, from this person's perspective.

Summary:
{summary}

###Response to evaluate:
{summary}

###Score Rubric:
{rubric}
\end{lstlisting}
\begin{tabular}{p{\textwidth}}
\bottomrule
\end{tabular}
\caption[Prompt template for LLM Judges.]{Prompt template for LLM Judges. Score rubrics in~\Cref{tab:ch-7-llm-judge-rubrics}. Source is also provided in prompt for consistency and relevance.}\label{tab:ch-7-prompt-llm-judge}
\end{table}

% ====================================================================
\begin{table*}[h]
\centering
\footnotesize
\begin{tabular}{p{.95\textwidth}}
\toprule
\textbf{Score Rubrics} \\
\midrule
\end{tabular}

    \begin{tabular}{p{.9\textwidth}}
        \texttt{Coherence} \\ \hdashfull 
    \end{tabular}
    \begin{lstlisting}[style=promptstyle]
Is the summary well-organized and easy to read?
[1] The summary is incoherent and very difficult to follow.
[2] The summary has poor organization with frequent disjointed transitions.
[3] The summary is somewhat organized but has noticeable structural issues.
[4] The summary is well-organized with only minor flow issues.
[5] The summary is excellently structured and flows naturally throughout.
    \end{lstlisting}

    \begin{tabular}{p{.9\textwidth}}
        \texttt{Consistency} \\ \hdashfull 
    \end{tabular}
    \begin{lstlisting}[style=promptstyle]
Is the summary factually consistent with the source document?
[1] The summary contains major factual errors or hallucinations not found in the source.
[2] The summary contains several factual inconsistencies with the source.
[3] The summary is mostly consistent but contains a few minor factual errors.
[4] The summary is largely consistent with only negligible inaccuracies.
[5] All information in the summary is fully consistent with the source document.
    \end{lstlisting}

    \begin{tabular}{p{.9\textwidth}}
        \texttt{Fluency} \\ \hdashfull 
    \end{tabular}
    \begin{lstlisting}[style=promptstyle]
Is the summary grammatically correct and well-written?
[1] The summary has many grammatical errors and is poorly written.
[2] The summary has frequent grammatical issues that hinder readability.
[3] The summary has occasional grammatical errors but is generally readable.
[4] The summary is well-written with only minor grammatical issues.
[5] The summary has flawless grammar and excellent writing quality.
    \end{lstlisting}

    \begin{tabular}{p{.9\textwidth}}
        \texttt{Relevance} \\ \hdashfull 
    \end{tabular}
    \begin{lstlisting}[style=promptstyle]
Does the summary capture the key information from the source?
[1] The summary is completely irrelevant and fails to capture any key information from the source document.
[2] The summary captures very little key information and misses most important points.
[3] The summary captures some key information but misses several important points.
[4] The summary captures most key information with only minor omissions.
[5] The summary captures all key points and essential information from the source.
    \end{lstlisting}

    \begin{tabular}{p{.9\textwidth}}
        \texttt{Informativeness} \\ \hdashfull 
    \end{tabular}
    \begin{lstlisting}[style=promptstyle]
How useful and informative is the summary to this specific person, given their background and information needs?
    
[1] The summary provides no useful information for this person's specific needs and background. It fails to address their domain or information requirements.
[2] The summary provides very little information relevant to this person's needs. It largely misses content that would be useful given their role and domain expertise.
[3] The summary provides some information relevant to this person's needs but lacks important domain-specific details or fails to address their stated information requirements.
[4] The summary is informative for this person's needs, covering most relevant details for their role and domain, with only minor gaps in addressing their information requirements.
[5] The summary is highly informative for this person's specific needs, thoroughly addressing their domain expertise and information requirements with relevant detail and appropriate context. 
    \end{lstlisting}
    
\begin{tabular}{p{\textwidth}}
\bottomrule
\end{tabular}
\caption[Score rubrics.]{Score rubrics for prompt in~\Cref{tab:ch-7-prompt-llm-judge}.}\label{tab:ch-7-llm-judge-rubrics}
\end{table*}

% ====================================================================

\begin{table}[h]
\centering
\footnotesize
\begin{tabular}{p{.9\columnwidth}}
\toprule
\textbf{Prompt: LLM Judge Annotator} \\
\midrule
\end{tabular}
\begin{lstlisting}[style=promptstyle]
###Task Description:
You are evaluating a summary from the perspective of a specific annotator. Consider their background and information needs when assessing quality.
1. Write a detailed feedback that assesses strictly whether the summary best addresses the annotator's query, not the fidelity or the style.
2. After writing a feedback, write a score that is an integer between 1 and 5. You should refer to the score rubric.
3. The output format should look as follows: "Feedback: (write a feedback for criteria) [RESULT] (an integer number between 1 and 5)"
4. Please do not generate any other opening, closing, or explanations.

###Introduction
Automatic summarization tools condense large amounts of text into shorter versions, highlighting the most important points.
But "important" is subjective - what matters to one reader might not matter to another.
This study explores what individual readers actually want from summaries of scientific papers, and how we can measure whether a summary was useful to them.
**We're focused on whether the summary answers your query, not the style or fidelity of the summary.**

###Annotator profile:
- Role: {role}
- Domain: {domain}
- Information needs: {info_needs}

###Query:
{query}

###Response to evaluate:
{summary}

###Score Rubric:
How well does the summary address the annotator's specific query, given their background and information needs?
        
[1] The summary does not address the annotator's query at all. It provides no information that would help answer what they asked about.
[2] The summary barely addresses the query, touching on it only tangentially or providing very little of the requested information.
[3] The summary partially addresses the query, providing some relevant information but missing important aspects of what was asked.
[4] The summary addresses the query well, covering most of what was asked with only minor gaps.
[5] The summary fully and directly addresses the annotator's query, providing the requested information clearly and completely.
\end{lstlisting}
\begin{tabular}{p{\textwidth}}
\bottomrule
\end{tabular}
\caption[Prompt template for LLM Judge Annotator.]{Prompt template for LLM Judge Annotator, in which the prompt reflects the instructions given to the human annotators.}\label{tab:ch-7-prompt-llm-judge-annotator}
\end{table}

% ====================================================================

\begin{table}[h]
\centering
\footnotesize
\begin{tabular}{p{.95\columnwidth}}
\toprule
\textbf{Prompt: Persona Precision} \\
\midrule
\end{tabular}

    \begin{tabular}{p{.9\columnwidth}}
        \texttt{Step: Extract} \\ \hdashfull 
    \end{tabular}
    \begin{lstlisting}[style=promptstyle]
Break down the following summary into a list of independent informational nuggets. Each nugget should be a single, self-contained piece of information -- a minimal, atomic statement that conveys one fact or claim.

## Summary
{summary}

## Instructions
List each informational nugget on its own line, prefixed with "- ". Output ONLY the list of nuggets, nothing else.

Example format:
- The study used a randomized controlled trial design.
- The sample included 500 participants from three hospitals.
- Treatment group showed 23% improvement over control.
    \end{lstlisting}

    \begin{tabular}{p{.9\columnwidth}}
    \texttt{Step: Relevance} \\ \hdashfull 
    \end{tabular}
    \begin{lstlisting}[style=promptstyle]
Determine whether the following informational nugget from a summary is relevant to the given persona and their query.

A nugget is "relevant" if it addresses the persona's information needs, relates to their domain of expertise, or helps answer their query. A nugget is "not relevant" if it is unrelated to what this persona would care about.

## Persona
- Role: {role}
- Domain: {domain}
- Information needs: {info_needs}

## Query
{query}

## Nugget
{nugget}

## Instructions
Is this nugget relevant to the persona's needs and query? Respond with ONLY "relevant" or "not relevant".
    \end{lstlisting}


\begin{tabular}{p{\columnwidth}}
\bottomrule
\end{tabular}
\caption{Prompts for Persona Precision metric.}\label{tab:ch-7-llm-judge-persona-precision}
\end{table}

% ====================================================================
\begin{table}[h]
\centering
\scriptsize
\setlength{\abovecaptionskip}{4pt}
\vspace{-4pt}
\begin{tabular}{p{.95\columnwidth}}
\toprule
\textbf{Prompt: Persona Recall} \\
\midrule
\end{tabular}

\noindent\texttt{Step: Extract}\dotfill\par
\vspace{-2pt}
\begin{lstlisting}[style=promptstyle,aboveskip=2pt,belowskip=2pt]
Break down the following summary into a list of independent informational nuggets. Each nugget should be a single, self-contained piece of information -- a minimal, atomic statement that conveys one fact or claim.

## Summary
{summary}

## Instructions
List each informational nugget on its own line, prefixed with "- ". Output ONLY the list of nuggets, nothing else.

Example format:
- The study used a randomized controlled trial design.
- The sample included 500 participants from three hospitals.
- Treatment group showed 23% improvement over control.
\end{lstlisting}
\vspace{-2pt}

\noindent\texttt{Step: Requirements}\dotfill\par
\vspace{-2pt}
\begin{lstlisting}[style=promptstyle,aboveskip=2pt,belowskip=2pt]
You are given a persona description and their query about a source document. Generate a list of informational nugget requirements -- specific pieces of information that a good summary SHOULD contain to serve this persona's needs.

Each requirement should be an atomic, specific piece of information that this persona would expect or need from a summary of the source document.

## Persona
- Role: {role}
- Domain: {domain}
- Information needs: {info_needs}

## Query
{query}

## Source Document
{source}

## Instructions
Based on the persona's background, information needs, and query, list the key informational nuggets that a good summary should include. Each nugget should be a single, specific piece of information from the source document.

List each requirement on its own line, prefixed with "- ". Output ONLY the list of requirements, nothing else.

Example format:
- The main methodology used in the study.
- Key quantitative results and effect sizes.
- Limitations that affect practical application.
\end{lstlisting}
\vspace{-2pt}

\noindent\texttt{Step: Coverage}\dotfill\par
\vspace{-2pt}
\begin{lstlisting}[style=promptstyle,aboveskip=2pt,belowskip=2pt]
Determine whether the following informational requirement is covered by any of the summary nuggets below.

A requirement is "covered" if any of the summary nuggets address the same information, even if phrased differently.

## Requirement
{requirement}

## Summary Nuggets
{nuggets}

## Instructions
Is the requirement addressed by any of the summary nuggets above? Respond with ONLY "covered" or "not covered".
\end{lstlisting}
\vspace{-2pt}

\begin{tabular}{p{\columnwidth}}
\bottomrule
\end{tabular}
\vspace{-6pt}
\caption{Prompts for Persona Recall metric.}\label{tab:ch-7-llm-judge-persona-recall}
\end{table}

% ====================================================================

\begin{table}[h]
\centering
\footnotesize
\begin{tabular}{p{.95\columnwidth}}
\toprule
\textbf{Prompt: Perturbation Prompts} \\
\midrule
\end{tabular}

    \begin{tabular}{p{.9\columnwidth}}
        \texttt{Test: Lengthen} \\ \hdashfull 
    \end{tabular}
    \begin{lstlisting}[style=promptstyle]
You are given a summary and the source document it was based on.

Rewrite the summary to be longer and more detailed. You MUST follow these rules:
- Do NOT add any new information that is not already present in or directly inferable from the original summary.
- Keep all original facts, claims, and findings intact.
- Expand by elaborating on existing points, adding transitional phrases, and using more descriptive language.
- The rewritten summary should be approximately 50% longer than the original.

Source document:
{source}

Original summary:
{summary}

Rewritten (longer) summary:
    \end{lstlisting}

    \begin{tabular}{p{.9\columnwidth}}
    \texttt{Test: Shorten} \\ \hdashfull 
    \end{tabular}
    \begin{lstlisting}[style=promptstyle]
You are given a summary and the source document it was based on.

Rewrite the summary to be shorter and more concise. You MUST follow these rules:
- Do NOT remove any key information, facts, claims, or findings from the original summary.
- Preserve all substantive content while using fewer words.
- Remove redundancy, simplify phrasing, and tighten sentence structure.
- The rewritten summary should be approximately 50% shorter than the original.

Source document:
{source}

Original summary:
{summary}

Rewritten (shorter) summary:
    \end{lstlisting}

    \begin{tabular}{p{.9\columnwidth}}
    \texttt{Test: Different Audience} \\ \hdashfull 
    \end{tabular}
    \begin{lstlisting}[style=promptstyle]
You are given a summary of a source document. The summary was originally written for: {original_audience}.

Rewrite this summary for a different target audience: {target_audience}.

You should:
- Adjust the language, terminology, and level of detail to be appropriate for the target audience.
- You may change which information is emphasized or how concepts are explained.
- Base your rewrite on the source document, not just the original summary.
- The rewrite should be roughly the same length as the original summary.

Source document:
{source}

Original summary (written for {original_audience}):
{summary}

Rewritten summary (for {target_audience}):
    \end{lstlisting}


\begin{tabular}{p{\columnwidth}}
\bottomrule
\end{tabular}
\caption[Prompts for the synthetic perturbations tests.]{Prompts for the synthetic perturbations tests, described in~\Cref{sec:ch-7-experimental-setup}.}\label{tab:ch-7-perturbation-prompts}
\end{table}